\section{Conclusiones}
\label{sec:conclusiones}

\subsection{Síntesis de los hallazgos}

El proyecto diseñó, implementó y validó estadísticamente un módulo de software
analítico basado en el Filtro de Kalman Lineal Discreto, concebido como capa de
consistencia física entre la telemetría degradada de un chancador primario y las
capas superiores de decisión de un gemelo digital.

Sobre la variable observada, el módulo redujo el error cuadrático medio del
volumen de carga en un \SI{89,57}{\percent} respecto de la lectura SCADA cruda,
reducción estadísticamente significativa según prueba t pareada
($t(49) = 135{,}24$; $p < 0{,}001$) y confirmada por contraste no paramétrico de
Wilcoxon ante el rechazo del supuesto de normalidad. Sobre la variable latente
—la tasa de llenado, que ningún instrumento mide— el módulo redujo el error en
un \SI{98,17}{\percent} respecto del único procedimiento disponible para un
SCADA convencional, produciendo información de estado que antes no existía en el
sistema.

El análisis de sensibilidad acreditó que el aporte del módulo crece con la
degradación del entorno, y que su desempeño está gobernado por la razón entre
las covarianzas de proceso y medición antes que por sus valores absolutos,
hallazgo que reduce sustantivamente la dificultad práctica de sintonización.

\subsection{Cumplimiento de los objetivos}

Los cuatro objetivos específicos se cumplieron, y cada uno generó el entregable
auditable comprometido en el plan del proyecto: matrices de covarianza
parametrizadas con fundamento físico trazable (OE1), un modelo en espacio de
estados con observabilidad verificada formalmente (OE2), un módulo de software
operativo, desacoplado y reproducible (OE3), y una validación estadística
completa con análisis de robustez (OE4). El objetivo general —diseñar y validar
mediante simulación estocástica el módulo, optimizando la observabilidad y
reduciendo el error cuadrático medio— se alcanzó dentro del alcance declarado.

\subsection{Aportes del estudio}

El aporte del trabajo no reside en el algoritmo, que es teoría establecida desde
\textcite{kalman1960}, sino en tres planos distintos.

En el plano \textbf{arquitectónico}, el estudio formula la estimación estocástica
de estados como componente de diseño de un gemelo digital minero, y no como
instrumento de control ni como etapa de preprocesamiento, atendiendo la capa que
las revisiones de \textcite{qu2023} y \textcite{ebad2026} identifican como
desatendida.

En el plano \textbf{empírico}, aporta evidencia cuantificada de que la
diferencia entre una sombra digital y un gemelo digital es medible, y de que la
capa analítica puede generar observabilidad allí donde la instrumentación no la
provee.

En el plano \textbf{metodológico}, entrega un artefacto reproducible byte a
byte, con registros de trazabilidad, verificación de supuestos estadísticos y
análisis de robustez paramétrica, en un área donde la validación efectiva ha sido
señalada como carencia.

\subsection{Recomendaciones y líneas de continuidad}

Se identifican cuatro líneas de continuidad, ordenadas por proximidad.

La primera y más inmediata es la \textbf{validación con telemetría real} de una
faena operativa, que permitiría contrastar el desempeño bajo discrepancia
genuina de modelo y determinar el nivel de madurez tecnológica efectivamente
alcanzable. Esta línea es condición necesaria para cualquier despliegue.

La segunda es la \textbf{extensión a regímenes no lineales}, incorporando los
estados transitorios de arranque y parada y los eventos de atascamiento que este
estudio excluyó, lo que requeriría migrar hacia variantes extendida o
\emph{unscented} del estimador, siguiendo el camino de \textcite{barzaga2025}.

La tercera es la \textbf{sintonización adaptativa} de las covarianzas, sugerida
directamente por el hallazgo sobre la razón $\mathbf{Q}/\mathbf{R}$: un
mecanismo que ajuste esa proporción en línea a partir de la secuencia de
innovaciones reduciría aún más la barrera de adopción industrial.

La cuarta es la \textbf{generalización del módulo como servicio transversal} del
gemelo digital, aplicando el mismo núcleo a los distintos subsistemas del activo
según la evidencia de extensibilidad reportada en la
sección~\ref{sec:extensibilidad}.

\subsection{Reflexión final}

La afirmación que sostiene este proyecto es que la inteligencia artificial
industrial requiere obligatoriamente limpiar matemáticamente sus datos antes de
tomar decisiones operacionales. Los resultados obtenidos permiten precisarla:
no se trata solo de limpiar, sino de \emph{reconstruir}. En un sistema
parcialmente observable, la capa de consistencia física no mejora la información
disponible, sino que produce información que la instrumentación no puede
entregar. Reconocer esa función es lo que justifica tratarla como componente
arquitectónico de primer orden y no como un accesorio del \emph{pipeline} de
datos.
